
\marginnote{beginning of learning/triangle.tex}


\subsection{Learning a Watson Distribution} 

\experiment{1.grammars/test\_watson/test\_watson.py}

In the following experiments, we select a random triangle (by using a
Gaussian in Bookstein coordinates). We then draw 20 samples from the
Watson distribution centered at this triangle (using 30 for the
concentration parameter of the Watson). We then reestimate the Watson
distribution from the samples.

This is a less noisy version of the learning task that EM faces when
refitting the midpoint distributions of a grammar from 20 samples.

\begin{figure}
\includegraphics[width=0.6in]{output/1.models/test_watson/watson_1_true.png}\\ 
\includegraphics[width=4in]{output/1.models/test_watson/watson_1_samples.png}\\ 
\includegraphics[width=0.6in]{output/1.models/test_watson/watson_1_est.png}
\caption{Experimenting with the Watson distribution, part 1. In the first row, the original triangle $T$. Subsequent rows are samples from Watson$(T,30.00)$. The final row is the mode of the estimated Watson distribution. The estimated concentration was $25.59$.}
\label{fig-watson-1}
\end{figure}

\begin{figure}
\includegraphics[width=0.6in]{output/1.models/test_watson/watson_2_true.png}\\ 
\includegraphics[width=4in]{output/1.models/test_watson/watson_2_samples.png}\\ 
\includegraphics[width=0.6in]{output/1.models/test_watson/watson_2_est.png}
\caption{Experimenting with the Watson distribution, part 2. In the first row, the original triangle $T$. Subsequent rows are samples from Watson$(T,30.00)$. The final row is the mode of the estimated Watson distribution. The estimated concentration was $36.92$.}
\label{fig-watson-2}
\end{figure}

\begin{figure}
\includegraphics[width=0.6in]{output/1.models/test_watson/watson_3_true.png}\\ 
\includegraphics[width=4in]{output/1.models/test_watson/watson_3_samples.png}\\ 
\includegraphics[width=0.6in]{output/1.models/test_watson/watson_3_est.png}
\caption{Experimenting with the Watson distribution, part 3. In the first row, the original triangle $T$. Subsequent rows are samples from Watson$(T,30.00)$. The final row is the mode of the estimated Watson distribution. The estimated concentration was $21.55$.}
\label{fig-watson-3}
\end{figure}

\begin{figure}
\includegraphics[width=0.6in]{output/1.models/test_watson/watson_4_true.png}\\ 
\includegraphics[width=4in]{output/1.models/test_watson/watson_4_samples.png}\\ 
\includegraphics[width=0.6in]{output/1.models/test_watson/watson_4_est.png}
\caption{Experimenting with the Watson distribution, part 4. In the first row, the original triangle $T$. Subsequent rows are samples from Watson$(T,30.00)$. The final row is the mode of the estimated Watson distribution. The estimated concentration was $92.59$.}
\label{fig-watson-4}
\end{figure}

\begin{figure}
\includegraphics[width=0.6in]{output/1.models/test_watson/watson_5_true.png}\\ 
\includegraphics[width=4in]{output/1.models/test_watson/watson_5_samples.png}\\ 
\includegraphics[width=0.6in]{output/1.models/test_watson/watson_5_est.png}
\caption{Experimenting with the Watson distribution, part 5. In the first row, the original triangle $T$. Subsequent rows are samples from Watson$(T,30.00)$. The final row is the mode of the estimated Watson distribution. The estimated concentration was $31.68$.}
\label{fig-watson-5}
\end{figure}



\subsection{Non-parametric models}


Since we are modeling the midpoint distributions as non-parametric
distributions, the learning algorithm is straightforward once we have
observations of each midpoint distribution. If we see midpoints $z_i$,
each with soft count $c_i$, our midpoint distribution will be
$$P(z) = \frac{1}{2\pi h^2\sum_i c_i} \sum_i c_i e^{-\frac{\| z - z_i\|^2}{2h^2}}$$

In vision applications where we match data to a model by minimizing a
cost function, it is often best to limit the cost of matching any one
piece of the model. In this case, we may want to limit the cost of
matching any two midpoints to some maximum value. We may therefore
also add a distribution that is uniform over the set of ``reasonable''
midpoints. If this set has area $A$, then our midpoint distribution
will be
$$P(z) = \frac{1}{2\pi h^2\sum_i c_i} \left[2\pi h^2 \cdot \frac{c_0}{A} + \sum_{i>0} c_i
  e^{-\frac{\| z - z_i\|^2}{2h^2}}\right].$$

We build our initial grammar by decomposing individual training
examples (see Section \ref{sec-single}), so our initial midpoint
distributions are simply Gaussians. Our initial Parzen window size
$h$, which is the standard deviation of the Gaussian, will be set
significantly higher than our post-training window size.

\marginnote{end of learning/triangle.tex}
